% Image credits for the photographs and AI illustrations of Book 4
% (University Chemistry, Year 3). Machine-readable license records live in
% images/book4/CREDITS.md; keep the two files in sync. The Book 4 writer
% owns this file: add one \item per photograph (in an itemize, after the
% first paragraph) as photographs land.
% Arabic edition: personal names are transliterated as in the captions, with
% the original spelling in parentheses (the legally required credit); user
% names, institutions, works and licences stay verbatim (as in the Arabic
% Book 2 credits page, image-credits-book2.ar.tex).
\chapter*{حقوق الصور}

رسوم هذا الكتاب أصلية. أما الصور الفوتوغرافية، حيث استُعملت، فمستعملة
بمقتضى رخصها الحرة، مع الشكر لأصحابها؛ وروابط المصادر الكاملة وروابط
الرخص لكل صورة مدرجة في
\texttt{images/book4/CREDITS.md} في مستودع الكتاب.
والرسوم التحذيرية للأخطار هي رسوم النظام المنسَّق عالميًا لتصنيف
المواد الكيميائية ووسمها، كما رُسمت لحزمة
\texttt{ghsystem} في \TeX\ Live.

\bigskip
وإلى جانب الصور الفوتوغرافية والرسوم الأصلية، بعض الأشكال
رسوم مولَّدة بالذكاء الاصطناعي، أُنتجت بأداة توليد الصور من OpenAI انطلاقًا من
أوامر كتبها محررو الكتاب، ورُوجعت من حيث الدقة الكيميائية قبل
إدراجها. والأمر الكامل لكل صورة مولَّدة
مدوَّن في \texttt{images/book4/ai/PROMPTS.md} في المستودع.
\bigskip
\noindent الصور الفوتوغرافية، حسب الفصل:
\begin{itemize}\small
  \item الفصل~1: إرفين شرودنغر (Erwin Schr\"odinger)، 1933، Nobel Foundation، ملك عام
        (Wikimedia Commons).
  \item الفصل~4: بلورات الثلج، تصوير ويلسون أ. بنتلي (Wilson A. Bentley) (نحو 1902)، ملك عام
        (Wikimedia Commons).
  \item الفصل~6: هوائيات راديوية على هضبة مرتفعة، ESO/ب.~تفرشي (B.~Tafreshi)
        (twanight.org)، CC BY 4.0 (Wikimedia Commons).
  \item الفصل~7: الفلوريت في ضوء النهار وتحت الأشعة فوق البنفسجية، تصوير ماشا
        ميلشينا (Masha Milshina)، CC BY 4.0 (Wikimedia Commons).
  \item الفصل~8: مختبر للرنين المغناطيسي النووي، تصوير Shandchem، CC BY 2.0 (Wikimedia
        Commons).
  \item الفصل~9: ويليام لورنس براغ (William Lawrence Bragg)، Nobel Foundation، ملك عام
        (Wikimedia Commons).
  \item الفصل~10: قبر لودفيغ بولتزمان، تصوير Feldkurat Katz،
        CC BY-SA 4.0 (Wikimedia Commons).
  \item الفصل~13: حلزونات بيلوسوف--جابوتينسكي، تصوير ستيفن موريس (Stephen Morris)، CC BY 2.0
        (Wikimedia Commons).
  \item الفصل~14: ثقب الأوزون فوق القارة القطبية الجنوبية، NASA، ملك عام (Wikimedia
        Commons).
  \item الفصل~16: صورة بالمجهر الإلكتروني لقرص خلوي من محوِّل حفزي، تصوير
        Galadrid، CC BY-SA 4.0 (Wikimedia Commons).
  \item الفصل~17: أشعة الشمس ومفعول تندال، تصوير كيفن هيغينز (Kevin c higgins)، CC BY-SA 4.0
        (Wikimedia Commons).
  \item الفصل~18: بلورات الياقوت الأحمر والزمرد، تصوير روبرت م. لافينسكي (Robert M. Lavinsky)، CC BY-SA 3.0
        (Wikimedia Commons).
  \item الفصل~20: بلورات الفيروسين، تصوير أندريا سيلا (Andrea Sella)، CC BY 4.0 (Wikimedia
        Commons).
  \item الفصل~22: سبيكة سيليكون أحادية البلورة، تصوير سيباستيان فالروت (Sebastian Wallroth)، ملك
        عام (Wikimedia Commons).
  \item الفصل~23: زهرة على هلام هوائي من السيليكا، NASA/JPL، ملك عام؛ و
        كأس ليكورغوس، تصوير ماري-لان نغوين (Marie-Lan Nguyen)، CC BY 2.5 (Wikimedia Commons).
  \item الفصل~24: سلطعون حدوة الحصان الأطلسي، تصوير Kaldari، CC0 (Wikimedia
        Commons).
  \item الفصل~27: أقحوان البيرثروم، تصوير Soramimi، CC BY-SA 4.0 (Wikimedia
        Commons).
  \item الفصل~30: لحاء طقسوس الباسيفيك، تصوير JOE BLOWE (theforestprimeval)،
        CC BY-SA 2.0 (Wikimedia Commons).
  \item الفصل~32: زهرة طحالب كما رآها القمر الاصطناعي Landsat 8، USGS وNASA،
        ملك عام.
  \item الفصل~33: مشعب مزدوج للتفريغ وللغاز الخامل، تصوير Polimerek، CC BY-SA
        3.0 (Wikimedia Commons).
\end{itemize}
\clearpage
